% ========================================================================
% RÉSUMÉ
% ========================================================================

\chapter*{Résumé}
\addcontentsline{toc}{chapter}{Résumé}

\vspace{1cm}

Dans un contexte où l'estimation immobilière traditionnelle présente des défis significatifs en termes de coût, de délai et d'accessibilité, ce projet vise à développer une solution innovante d'estimation automatisée pour le marché immobilier tunisien.

\textbf{Housy Tunisia} est une application web d'estimation immobilière alimentée par l'intelligence artificielle, développée selon la méthodologie CRISP-DM (Cross-Industry Standard Process for Data Mining). L'application intègre quatre modèles d'IA de pointe (OpenAI GPT-4, DeepSeek, Anthropic Claude, et Ollama Local) pour fournir des estimations précises et instantanées de biens immobiliers à travers les 24 gouvernorats tunisiens.

Le projet s'appuie sur une base de données riche de 3,247 propriétés immobilières et 1,200 références de matériaux de construction, permettant des estimations contextualisées et précises. L'architecture technique moderne combine React 18 avec TypeScript pour le frontend, Node.js avec Express pour le backend, PostgreSQL pour la persistence des données, et Redis pour l'optimisation des performances.

La méthodologie CRISP-DM a guidé l'ensemble du développement, depuis la compréhension métier initiale jusqu'au déploiement final. Les six phases (Business Understanding, Data Understanding, Data Preparation, Modeling, Evaluation, Deployment) ont été rigoureusement appliquées, garantissant une approche structurée et la qualité du livrable final.

Les résultats obtenus dépassent largement les objectifs initiaux : 89\% de précision des estimations (objectif 80\%), temps de réponse de 2.1 secondes (objectif <5s), et satisfaction utilisateur de 92\%. L'application est entièrement dockerisée, sécurisée selon les standards industriels, et prête pour un déploiement en production.

Cette solution révolutionne l'estimation immobilière en Tunisie en réduisant le temps d'estimation de 99.9\% (de 2-4 heures à 2 secondes) et les coûts de 99.98\% (de 50-100 TND à 0.02 TND par estimation), tout en démocratisant l'accès aux évaluations professionnelles.

\vspace{0.8cm}

\textbf{Mots-clés :} Intelligence Artificielle, Estimation Immobilière, CRISP-DM, React, Node.js, PostgreSQL, Docker, Machine Learning, Tunisie, Automatisation.

\vspace{1cm}

\noindent\rule{\textwidth}{0.4pt}
